\documentclass[../../main.tex]{subfiles}

\begin{document}

\section{Lebesgue 可测集的外逼近和内逼近}\label{section:ch02-内外逼近}

我们现在转向基于闭集的内逼近与基于开集的外逼近的单个集合的可测性的特征, 这些特征提供了可测性的交替视角, 是我们以后研究关于可测与可积函数的逼近性质的必要工具.

可测集具有以下分割性质: 若 $A$ 是包含于 $B$ 的有限外测度的可测集, 则
\begin{align}
m^{*}(B\sim A)=m^{*}(B)-m^{*}(A).\label{align:ch02-分割性质}
\end{align}
事实上, 由 $A$ 的可测性,
\begin{align*}
m^{*}(B)=m^{*}(B\cap A)+m^{*}(B\cap A^{c})=m^{*}(A)+m^{*}(B\sim A).
\end{align*}
因此, 由于 $m^{*}(A)<\infty$, 我们有\eqref{align:ch02-分割性质}.

\begin{theorem}\label{theorem:ch02-可测等价}
令 $E$ 为任何实数集. 则下列四个断言都与 $E$ 的可测性等价.

\textbf{(用开集和 $G_{\delta}$ 集的外逼近)}
\begin{enumerate}[(i)]
\item 对每个 $\varepsilon>0$, 存在包含 $E$ 的开集 $\mathcal{O}$ 使得 $m^{*}(\mathcal{O}\sim E)<\varepsilon$.
\item 存在包含 $E$ 的 $G_{\delta}$ 集 $G$ 使得 $m^{*}(G\sim E)=0$.
\end{enumerate}

\textbf{(用闭集和 $F_{\sigma}$ 集的内逼近)}
\begin{enumerate}[(iii)]
\item 对每个 $\varepsilon>0$, 存在包含于 $E$ 的闭集 $F$ 使得 $m^{*}(E\sim F)<\varepsilon$.
\item 存在包含于 $E$ 的 $F_{\sigma}$ 集 $F$ 使得 $m^{*}(E\sim F)=0$.
\end{enumerate}
\end{theorem}

\begin{proof}
我们证明 $E$ 的可测性与两个外逼近性质 (i) 和 (ii) 的等价性. 证明的剩余部分从 De Morgan 等式以及以下观察得到: 一个集合是可测的当且仅当它的补是可测的, 一个集合是开的当且仅当它的补是闭的, 一个集合是 $F_{\sigma}$ 当且仅当它的补是 $G_{\delta}$.

假设 $E$ 是可测的. 令 $\varepsilon>0$. 首先考虑 $m^{*}(E)<\infty$ 的情形. 根据外测度的定义, 存在开区间的可数族 $\{I_{k}\}_{k=1}^{\infty}$ 覆盖 $E$ 且使得
\begin{align*}
\sum_{k=1}^{\infty}\ell(I_{k})<m^{*}(E)+\varepsilon.
\end{align*}
定义 $\mathcal{O}=\bigcup_{k=1}^{\infty}I_{k}$. 则 $\mathcal{O}$ 是包含 $E$ 的开集. 根据 $\mathcal{O}$ 的外测度的定义,
\begin{align*}
m^{*}(\mathcal{O})\leqslant\sum_{k=1}^{\infty}\ell(I_{k})<m^{*}(E)+\varepsilon.
\end{align*}
因此,
\begin{align*}
m^{*}(\mathcal{O})-m^{*}(E)<\varepsilon.
\end{align*}
然而, $E$ 是可测的且有有限外测度. 因此, 根据上面注意到的可测集的分割性质,
\begin{align*}
m^{*}(\mathcal{O}\sim E)=m^{*}(\mathcal{O})-m^{*}(E)<\varepsilon.
\end{align*}

现在考虑 $m^{*}(E)=\infty$ 的情形. 则 $E$ 可表示为可测集的可数族 $\{E_{k}\}_{k=1}^{\infty}$ 的不交并, 该族的每个集合具有有限外测度. 根据有限外测度的情形, 对每个指标 $k$, 存在包含 $E_{k}$ 的开集 $\mathcal{O}_{k}$ 使得 $m^{*}(\mathcal{O}_{k}\sim E_{k})<\varepsilon/2^{k}$. 集合 $\mathcal{O}=\bigcup_{k=1}^{\infty}\mathcal{O}_{k}$ 是开的, 它包含 $E$ 且
\begin{align*}
\mathcal{O}\sim E=\bigcup_{k=1}^{\infty}\mathcal{O}_{k}\sim E\subseteq\bigcup_{k=1}^{\infty}[\mathcal{O}_{k}\sim E_{k}].
\end{align*}
因此,
\begin{align*}
m^{*}(\mathcal{O}\sim E)\leqslant\sum_{k=1}^{\infty}m^{*}(\mathcal{O}_{k}\sim E_{k})<\sum_{k=1}^{\infty}\varepsilon/2^{k}=\varepsilon.
\end{align*}
于是性质 (i) 对 $E$ 成立.

现在假设性质 (i) 对 $E$ 成立. 对每个自然数 $k$, 选取包含 $E$ 的开集 $\mathcal{O}_{k}$ 使得 $m^{*}(\mathcal{O}_{k}\sim E)<1/k$. 定义 $G=\bigcap_{k=1}^{\infty}\mathcal{O}_{k}$. 则 $G$ 是一个包含 $E$ 的 $G_{\delta}$ 集. 此外, 由于对每个 $k$, $G\sim E\subseteq\mathcal{O}_{k}\sim E$, 根据外测度的单调性,
\begin{align*}
m^{*}(G\sim E)\leqslant m^{*}(\mathcal{O}_{k}\sim E)<1/k.
\end{align*}
因此 $m^{*}(G\sim E)=0$, 从而 (ii) 成立. 现在假设性质 (ii) 对 $E$ 成立. 由于测度为零的集合是可测的, 且 $G_{\delta}$ 集也如此, 可测集是一个代数, 集合
\begin{align*}
E=G\cap[G\sim E]^{c}
\end{align*}
是可测的.
\end{proof}

有限外测度的可测集的以下性质断言这样的集合 "接近" 于有限个开区间的不交并.

\begin{theorem}\label{theorem:ch02-有限区间逼近}
令 $E$ 为外测度有限的可测集. 则对每个 $\varepsilon>0$, 存在开区间的有限不交族 $\{I_{k}\}_{k=1}^{n}$, 使得若 $\mathcal{O}=\bigcup_{k=1}^{n}I_{k}$, 则
\begin{align*}
m^{*}(E\sim\mathcal{O})+m^{*}(\mathcal{O}\sim E)<\varepsilon.
\end{align*}
\end{theorem}

\begin{proof}
根据\cref{theorem:ch02-可测等价}的断言 (i), 存在开集 $\mathcal{U}$ 使得
\begin{align}
E\subseteq\mathcal{U}\quad\text{且}\quad m^{*}(\mathcal{U}\sim E)<\varepsilon/2.\label{align:ch02-逼近U}
\end{align}
由于 $E$ 是可测的且具有有限的外测度, 我们从外测度的分割性质推断出 $\mathcal{U}$ 也有有限的外测度. 每个实数的开集是开区间的可数族的不交并. 令 $\mathcal{U}$ 为开区间的可数不交族 $\{I_{k}\}_{k=1}^{\infty}$ 的并. 每个区间是可测的且它的外测度是它的长度. 因此, 根据\cref{proposition:ch02-有限不交可加}和外测度的单调性, 对每个自然数 $n$,
\begin{align*}
\sum_{k=1}^{n}\ell(I_{k})=m^{*}\left(\bigcup_{k=1}^{n}I_{k}\right)\leqslant m^{*}(\mathcal{U})<\infty.
\end{align*}
该不等式的右边与 $n$ 无关. 因此,
\begin{align*}
\sum_{k=1}^{\infty}\ell(I_{k})<\infty.
\end{align*}
选取自然数 $n$ 使得
\begin{align*}
\sum_{k=n+1}^{\infty}\ell(I_{k})<\varepsilon/2.
\end{align*}
定义 $\mathcal{O}=\bigcup_{k=1}^{n}I_{k}$. 由于 $\mathcal{O}\sim E\subseteq\mathcal{U}\sim E$, 根据外测度的单调性和\eqref{align:ch02-逼近U},
\begin{align*}
m^{*}(\mathcal{O}\sim E)\leqslant m^{*}(\mathcal{U}\sim E)<\varepsilon/2.
\end{align*}
另一方面, 由于 $E\subseteq\mathcal{U}$,
\begin{align*}
E\sim\mathcal{O}\subseteq\mathcal{U}\sim\mathcal{O}=\bigcup_{k=n+1}^{\infty}I_{k}.
\end{align*}
因此根据外测度的定义,
\begin{align*}
m^{*}(E\sim\mathcal{O})\leqslant\sum_{k=n+1}^{\infty}\ell(I_{k})<\varepsilon/2.
\end{align*}
因此,
\begin{align*}
m^{*}(\mathcal{O}\sim E)+m^{*}(E\sim\mathcal{O})<\varepsilon.
\end{align*}
\end{proof}

\begin{remark}
关于\cref{theorem:ch02-可测等价}的断言 (i) 的说明是妥当的. 根据外测度的定义, 对任何有界集 $E$（无论是否可测）和任何 $\varepsilon>0$, 存在开集 $\mathcal{O}$ 使得 $E\subseteq\mathcal{O}$ 且 $m^{*}(\mathcal{O})<m^{*}(E)+\varepsilon$, 因此 $m^{*}(\mathcal{O})-m^{*}(E)<\varepsilon$. 这不蕴涵着 $m^{*}(\mathcal{O}\sim E)<\varepsilon$, 因为分割性质
\begin{align*}
m^{*}(\mathcal{O}\sim E)=m^{*}(\mathcal{O})-m^{*}(E)
\end{align*}
不成立, 除非 $E$ 是可测的（见习题 19）.
\end{remark}

\begin{exercise}\label{exercise:ch02-16}
通过证明可测性等价于 (iii) 且等价于 (iv) 完成\cref{theorem:ch02-可测等价}的证明.
\end{exercise}

\begin{exercise}\label{exercise:ch02-17}
证明集合 $E$ 是可测的当且仅当对每个 $\varepsilon>0$, 存在闭集 $F$ 和开集 $\mathcal{O}$ 使得 $F\subseteq E\subseteq\mathcal{O}$ 且 $m^{*}(\mathcal{O}\sim F)<\varepsilon$.
\end{exercise}

\begin{exercise}\label{exercise:ch02-18}
令 $E$ 有有限外测度. 证明存在 $F_{\sigma}$ 集 $F$ 和 $G_{\delta}$ 集 $G$ 使得 $F\subseteq E\subseteq G$ 且 $m^{*}(F)=m^{*}(E)=m^{*}(G)$.
\end{exercise}

\begin{exercise}\label{exercise:ch02-19}
令 $E$ 有有限外测度. 证明: 若 $E$ 不可测, 则存在包含 $E$ 的开集 $\mathcal{O}$, 它具有有限外测度使得
\begin{align*}
m^{*}(\mathcal{O}\sim E)>m^{*}(\mathcal{O})-m^{*}(E).
\end{align*}
\end{exercise}

\begin{exercise}\label{exercise:ch02-20}
(Lebesgue) 令 $E$ 有有限外测度. 证明 $E$ 是可测的当且仅当对每个开的有界区间 $(a,b)$,
\begin{align*}
b-a=m^{*}((a,b)\cap E)+m^{*}((a,b)\sim E).
\end{align*}
\end{exercise}

\begin{exercise}\label{exercise:ch02-21}
用\cref{theorem:ch02-可测等价}的性质 (ii) 作为可测集的原始定义且证明两个可测集的并是可测的. 接着对性质 (iv) 做同样的事.
\end{exercise}

\begin{exercise}\label{exercise:ch02-22}
对任何集合 $A$, 定义 $m^{**}(A)\in[0,\infty]$ 为
\begin{align*}
m^{**}(A)=\inf\{m^{*}(\mathcal{O})\mid\mathcal{O}\supseteq A,\ \mathcal{O}\text{ 是开的}\}.
\end{align*}
这个集函数 $m^{**}$ 与外测度 $m^{*}$ 有何联系?
\end{exercise}

\begin{exercise}\label{exercise:ch02-23}
对任何集合 $A$, 定义 $m^{***}(A)\in[0,\infty]$ 为
\begin{align*}
m^{***}(A)=\sup\{m^{*}(F)\mid F\subseteq A,\ F\text{ 是闭的}\}.
\end{align*}
这个集函数 $m^{***}$ 与外测度 $m^{*}$ 有何联系?
\end{exercise}

\end{document}